\input{../tex-inputs/latex-leseansicht-vorspann}

\section[Arthur Schnitzler an Gustav Schwarzkopf, XXXX]{L04493 Arthur Schnitzler an Gustav Schwarzkopf, XXXX}
\nopagebreak\mylabel{L04493v}
\rehead{ }\normalsize\beginnumbering\briefempfaengerindex{Schwarzkopf, Gustav@\textsc{Schwarzkopf, Gustav}!zzzSchnitzler, Arthur@\emph{von Arthur Schnitzler}!1923-09-051@{5. 9. 1923}|(be}
\toendnotes[C]{\smallbreak\pagebreak[2]}
\correspDesc{Versand  durch Arthur Schnitzler am 5. 9. 1923 in Scuol
\newline{}Erhalt  durch Gustav Schwarzkopf im Zeitraum [6. 9. 1923
                  – 10. 9. 1923?] in Wien}\toendnotes[C]{\smallbreak}
\Standort{DLA, A:Schnitzler, HS.1985.1.1897.}
\physDesc{Bildpostkarte, 147 Zeichen
\newline{}Handschrift: Bleistift, deutsche Kurrent
\newline{}Versand: Stempel: »\nobreak{}\oindex{Scuol@\textbf{Scuol}|pwk}Schuls, 2. IX. 23, 12\nobreak{}«.  }\toendnotes[C]{\smallbreak}\pstart{}{\pb}\textsc{Oesterreich}\oindex{Österreich@\textbf{Österreich}|pw}\pend{}\pstart{}Hrn Guſtav Schwarzkopf\pend{}\pstart{}\textsc{Wien I}\oindex{I., Innere Stadt@\textbf{I., Innere Stadt}|pw}\pend{}\pstart{}\textsc{Tiefer Graben 17}\oindex{Wien@\textbf{Wien}!I., Innere Stadt@\textbf{I., Innere Stadt}!Tiefer Graben 17@\textbf{Tiefer Graben 17}|pw}\pend{}{\bigskip}
\pstart
           \noindent{}{\pb}\textcolor{gray}{\textbf{25081 Schuls. Hotel
                     Engadinerhof\oindex{Hotel Engadinerhof Tarasp Scuol@\textbf{Hotel Engadinerhof Tarasp Scuol}|pw}}}\pend
           \vspace{1em}
\pstart
           \raggedleft{}{\pb}Schuls\oindex{Scuol@\textbf{Scuol}|pw}, \label{K_L04493-1v}\edtext{5. 9. 23}{\lemma{\textnormal{\emph{5. 9. 23}}}\Cendnote{\textnormal{Der Poststempel datiert die Karte anders als Schnitzler auf den
                        2. 9. 1923, laut \emph{Tagebuch}\pwindex{Schnitzler, Arthur 15. 5. 1862 Wien – 21. 10. 1931 ebd.@\textsc{Schnitzler, Arthur} (15. 5. 1862 Wien – 21. 10. 1931 ebd.)!Tagebuch@\strich\emph{Tagebuch}|pwk} kam Schnitzler aber
                        erst am 3. 9. 1923 in Schuls\oindex{Scuol@\textbf{Scuol}|pwk}
                     an, vgl. A. S.: \emph{Tagebuch}, 3. 9. 1923. Ob der Irrtum bei Schnitzler lag oder die Post ein veraltetes Stempeldatum nutzte, ist nicht ermittelt.}}}\label{K_L04493-1}.\pend
           \vspace{0.5em}
\pstart
           Herzliche Grüße, auch an Ihren Bruder\pwindex{Schwarzkopf, Max 12.\,6.\,1857 Wien – 14.\,4.\,1928 ebd.@\textsc{Schwarzkopf, Max} (12.\,6.\,1857 Wien – 14.\,4.\,1928 ebd.)|pwv}; und auf baldiges Wiederſehen.\pend
           
\pstart
           Ihr{\\[\baselineskip]}\spacefill\mbox{Arthur}\pend
           \leftskip=0em{}\selectlanguage{ngerman}\endnumbering\briefempfaengerindex{Schwarzkopf, Gustav@\textsc{Schwarzkopf, Gustav}!zzzSchnitzler, Arthur@\emph{von Arthur Schnitzler}!1923-09-051@{5. 9. 1923}|)be}\mylabel{L04493h}
\begin{anhang}
\end{anhang}\newcommand{\dateiname}{L04493}\newcommand{\titel}{Arthur Schnitzler an Gustav Schwarzkopf, XXXX}\newcommand{\editorInnen}{Herausgegeben von Jahnke, SelmaMüller, Martin Anton}\input{../tex-inputs/latex-leseansicht-abspann}
